\documentclass[11pt]{article}
\usepackage[margin=1in]{geometry}
\usepackage{amsmath,amssymb,amsthm,hyperref}
\newtheorem{proposition}{Proposition}
\title{Wages with Additive Activities and Mobile Capital}
\author{Exploratory note for Oliver Pardo}
\date{22 November 2026}
\begin{document}
\maketitle
I study two activities producing the same final good, $Y=Y_S+Y_C$.
Workers are attached to their activity, $L_j=\ell_j N$, with fixed
$\ell_j>0$ and $\ell_S+\ell_C=1$. Both use the same labor-augmenting
technology $\dot A/A=\gamma$. Thus $H_j=AL_j$ grows at $n+\gamma$.
I hold AI efficiency $B$ fixed and begin with competitive AI supply,
at common service price $p=1/B$. There is no research in this benchmark.
The results about a growing equilibrium below are conditional on existence
and the explicitly stated limiting properties.

\section{Common capital allocation and wages}
The technologies are
\[
 Y_j=K_j^\alpha Z_j^{1-\alpha},\qquad
 Z_j=[\omega_{Lj}H_j^{r_j}+\omega_{Xj}X_j^{r_j}]^{1/r_j},\qquad
 r_j=\frac{\sigma_j-1}{\sigma_j},\quad 0<\sigma_C<1<\sigma_S.
\]
Each activity's positive weights sum to one. Capital is freely mobile.
Writing $Z=Z_S+Z_C$, its efficient allocation satisfies
$K_j/Z_j=K/Z$, and final production is $Y=K^\alpha Z^{1-\alpha}$.
Let
\begin{equation}
 Q=(1-\alpha)(K/Z)^\alpha.
 \label{eq:two-Q}
\end{equation}
This is the marginal value of either composite $Z_j$, not a second final-good
price. At an interior AI allocation,
\begin{equation}
 p=Q\omega_{Xj}(Z_j/X_j)^{1/\sigma_j},\qquad
 w_j=A Q\omega_{Lj}(Z_j/H_j)^{1/\sigma_j}.
 \label{eq:two-wages}
\end{equation}
These equations follow by differentiating final production after the optimal
capital allocation. They retain the envelope effect of reallocating capital.

\section{A conditional long-run wage ranking}
Define the substituting activity's AI-only productivity
$c_S=\omega_{XS}^{\sigma_S/(\sigma_S-1)}<1$.

\begin{proposition}[Wages in a regular AI-dominated equilibrium]
Suppose a competitive equilibrium has $X_S/H_S\to\infty$,
$K/Z\to\kappa>0$, and positive limiting $C/K$. Suppose also that the
limiting average output-per-worker growth $g$ exceeds $\gamma$.
Then the following average growth rates hold:
\begin{equation}
 g_{w_S}=\gamma+\frac{g-\gamma}{\sigma_S},\qquad
 g_{w_C}=\gamma.
 \label{eq:two-growth}
\end{equation}
Thus $\gamma<g_{w_S}<g$, $w_S/w_C\to\infty$, both labor-income
shares vanish, and $Y_C/Y\to0$. The complementary activity remains
active at all sufficiently late dates. Convergence of instantaneous growth
rates is not asserted without additional regularity.
\end{proposition}
\begin{proof}
As $X_S/H_S\to\infty$, $Z_S/X_S\to c_S$ and
$\partial Z_S/\partial X_S\to c_S$. The AI price condition yields
$Q\to Q_*=p/c_S$. In the complementary activity, write $v_C=X_C/H_C$.
Its right marginal composite product at $v_C=0$ is
$d_C=\omega_{XC}^{\sigma_C/(\sigma_C-1)}>1$.
For $v_C>0$, that marginal product decreases strictly from $d_C$ to zero.
Since $p/Q_*=c_S<1<d_C$, the optimal complementary activity is eventually
interior, and its ratio $v_C$ converges to a unique finite positive value.
Consequently $Z_C/H_C$, $K_C/H_C$, and $w_C/A$ converge to positive constants.
This establishes $g_{w_C}=\gamma$.

Since $Z_C$ grows at average rate $n+\gamma$ while $Z_S/H_S\to\infty$,
$Z_C/Z_S\to0$. Hence $Y_C/Y=Z_C/Z\to0$. The positive limit of $K/Z$
implies that $Z$, $Y$ and $K$ have the same average growth rate $n+g$.
Taking logarithms of the second expression in~\eqref{eq:two-wages}
for $j=S$ proves the first expression in~\eqref{eq:two-growth}.
The ratio $w_S/w_C$ therefore grows at average rate
$(g-\gamma)/\sigma_S>0$. Finally, $L_j/N=\ell_j$ is fixed and both wage
growth rates are below $g$. Thus $w_j L_j/Y\to0$ for each activity.
\end{proof}

The normalized complementary wage has the explicit positive limit
\begin{equation}
 \frac{w_C}{A}\longrightarrow
 \left[
 \frac{Q_*^{1-\sigma_C}-\omega_{XC}^{\sigma_C}p^{1-\sigma_C}}
 {\omega_{LC}^{\sigma_C}}
 \right]^{1/(1-\sigma_C)}.
 \label{eq:two-C-level}
\end{equation}
To obtain it, let $v_j=w_j/A$ denote the price of effective labor.
The CES unit-cost equality is
$Q^{1-\sigma_j}=\omega_{Lj}^{\sigma_j}v_j^{1-\sigma_j}
+\omega_{Xj}^{\sigma_j}p^{1-\sigma_j}$.
Solving for $v_C$ and taking $Q\to Q_*$ gives~\eqref{eq:two-C-level}.

Under the stated $C/K$ limit, the household Euler equation gives the
conditional candidate
\[
 g=\alpha[B(1-\alpha)c_S]^{(1-\alpha)/\alpha}-\delta-\rho.
\]
For the regime described above this candidate must exceed $\gamma$.
This expression does not establish existence or convergence from a given state.

\section{Interpretation and scope}
The substituting activity expands with capital and AI, and its remaining
workers become more productive. The complementary activity expands only
at the scale of its effective labor endowment in this limit. Its absolute
output and wages can grow even as its share of the economy disappears.
Neither this proof nor the conditional return formula establishes that an
unanticipated improvement in $B$ benefits every worker during the transition.
They compare long-run rates within a specified equilibrium, not welfare or
income levels across different technologies.

If both activities instead have finite positive stationary $K_j/H_j$ and
$X_j/H_j$, wages in both grow at $\gamma$. Existence of such an allocation
requires checking the common prices and household conditions.

With identical perfectly mobile workers, interior labor allocation requires
$w_S=w_C$. The divergent wage ratio in the proposition is therefore
incompatible with fixed positive worker shares in both activities under that
alternative mobility assumption. One must solve endogenous labor allocation
rather than apply~\eqref{eq:two-growth} to those workers.

With a monopoly charging a common price, the wage formulas still hold
conditional on a positive finite limiting price and the same real-allocation
limits. A discriminating monopoly can charge different prices across
activities. Neither monopoly equilibrium is solved here.
\end{document}
